\documentclass[a4paper,11pt,fleqn]{article}%fleqn pour aligner les equations à gauche
\usepackage{/home/rweye/Nextcloud/Documents/Overleaf/Styles/Cours}


\newcommand{\chnum}{Chapitre 2}
\newcommand{\chtitle}{Analyse d'un système}
\newcommand{\dtype}{Fiche Révision}

\begin{document}
\maketitle

\section*{Ce qu'il faut connaitre :}
\begin{multicols}{3}
\subsection*{Le vocabulaire :}
\begin{itemize}[label=$\square$, itemsep=0pt]
    \item composition d'un système
    \item titrage
    \item réaction support
    \item espèces titrantes et titrées
    \item équivalence
    \item proportions st\oe chiométriques
    \item réactif limitant, réactif en excès
    \end{itemize}
\subsection*{Les notions :}
    \begin{itemize}[label=$\square$, itemsep=0pt]
    \item spectre d'absorption
    \item couleurs complémentaires
    \item spectre infrarouge
    \item Principaux groupes caractéristiques et familles chimiques associées
    \end{itemize}
\subsection*{Les relations :}
    \begin{itemize}[label=$\square$, itemsep=0pt]
        \item Loi de Beer-Lambert
        \item Loi de Kohlrausch
        \item expression de la conductivité
    \end{itemize}
\subsection*{Les grandeurs :}
    \begin{itemize}[label=$\square$, itemsep=0pt]
        \item longueur d'onde
        \item concentration d'une solution
        \item concentration d'une espèce en solution
        \item absorbance
        \item conductivité
        \item conductance
        \item concentration
        \item concentration en masse
        \item titre massique
        \item masse volumique
        \item densité
        \item pH
    \end{itemize}
\end{multicols}

\section*{Ce qu'il faut savoir faire :}
\begin{table}[h!]
    \begin{tabular}{|m{.84\linewidth}|m{.04\linewidth}|m{.08\linewidth}|}
        \hline
        \textbf{La compétence} & \textbf{Où ?} & \textbf{Maîtrise ?} \\
        \hline
        Exploiter, à partir de données tabulées, un spectre d’absorption infrarouge ou UV-visible pour identifier un groupe caractéristique ou une espèce chimique. & & \\
        \hline
        Exploiter la loi de Beer-Lambert, la loi de Kohlrausch ou l’équation d’état des gaz parfaits pour déterminer une concentration en quantité de matière. Citer les domaines de validité de ces relations.& & \\
        \hline
        Mesurer une conductance et tracer une courbe d’étalonnage pour déterminer une concentration. & & \\
        \hline
        Réaliser une solution de concentration donnée en soluté, apportée à partir d’une solution de titre massique et de densité fournie. & & \\
        Établir la composition du système après ajout d’un volume de solution titrante, la transformation étant considérée comme totale. & & \\
        \hline
        Exploiter un titrage pour déterminer une quantité de matière, une concentration ou une masse & & \\
        \hline
        Dans le cas d’un titrage avec suivi par conductimétrie, justifier qualitativement l’évolution de la pente de la courbe à l’aide de données sur les conductivités ioniques molaires& & \\
        \hline
        Mettre en œuvre un suivi pH-métrique d’un titrage ayant pour support une réaction acide-base.& & \\
        \hline
        Mettre en œuvre le suivi conductimétrique d’un titrage.& & \\
        \hline
        Capacité numérique : représenter, à l’aide d’un langage de programmation, l’évolution des quantités de matière des espèces en fonction du volume de solution titrante versé. & & \\
        \hline
        Distinguer dosage et titrage& & \\
        \hline
    \end{tabular}
\end{table}

\end{document}
